\documentclass[10pt,a4paper]{amsart}
\usepackage[margin=19mm]{geometry}
\usepackage{amsmath,amssymb,mathtools}
\usepackage[scr=rsfs,cal=euler]{mathalfa}
\usepackage{xfrac}
\usepackage[unicode,hidelinks,bookmarks=false]{hyperref}
\newcommand{\ie}{i.e.\ }
\newcommand{\cf}{cf.\ }
\newcommand{\resp}{respectively\ }
\newcommand{\Subsection}{\subsection}
\providecommand{\nobreakdash}{\nobreak\hbox{-}}
\setlength{\emergencystretch}{2em}
\multlinegap0pt
\usepackage{bm}
\usepackage{bbm}
\usepackage{dsfont}
\usepackage{xspace}
\usepackage{cancel}
\usepackage{mathtools}
\usepackage{amsthm}
\makeatletter
\@ifundefined{lemma}{\newtheorem{lemma}{Lemma}}{}
\makeatother
\makeatletter
\expandafter\ifx\csname manualtheorem\endcsname\relax
\newenvironment{manualtheorem}[1]{%
\renewcommand\themanualtheoreminner{#1}%
  \manualtheoreminner
}{\endmanualtheoreminner}
\fi
\newcommand{\norm}[1]{\left \lVert #1 \right \rVert}
\expandafter\ifx\csname nscenter\endcsname\relax
\newenvironment{nscenter}
{\parskip=10pt\par\nopagebreak\centering}
{\par\noindent}
\fi
\makeatother
\title[Symmetric power of higher dimensional varieties]{Symmetric power of higher dimensional varieties}
\author{Ashima Bansal, Supravat Sarkar, Shivam Vats}
\date{Source: arXiv:2508.12654v1 (CC BY 4.0)}
\begin{document}
\maketitle

\noindent\emph{Source and licence.} Symmetric power of higher dimensional varieties. Version arXiv:2508.12654v1 (CC BY 4.0), distributed under the Creative Commons Attribution 4.0 International licence, as recorded in the arXiv licence metadata for this exact version and shown on the abstract page https://arxiv.org/abs/2508.12654v1. Full rights notice: licenses/arxiv-2508.12654-CC-BY-4.0.txt.\par\smallskip

\noindent\textit{Editorial scope.} Counted unit: the Lemma whose source label is \texttt{local structure} in the article (source0.tex lines 193--200, source label \texttt{local structure}), delivered together with the article's own complete original proof of it (source0.tex lines 201--254, one \texttt{proof} environment of 54 lines). No mathematics is written, repaired or re-derived: statement and proof are the article's own text line for line. Required context retained uncounted: none. Every definition, display and label of those lines is retained unchanged. Omitted from this package and not redistributed: 1--192, 255--416. Nothing inside a retained excerpt is omitted or altered: every retained body is delivered exactly as the article's lines are written, so no omission receipt of any kind accompanies this package. Citations: the retained excerpts carry no citation command at all, so no bibliography entry of the article is transcribed and no reference is redistributed. Reference adaptation: none was needed and none was made; every reference inside the delivered unit resolves inside it, so no unresolved reference or citation is produced. Printed slips kept literal: no printed word, symbol, hypothesis or number of the counted statement or of its proof was altered. Layout and class: the article's own class is \texttt{amsart} with packages including bm, graphicx, adjustbox, amsthm, amssymb, amsmath, fontenc, inputenc, hyperref, xcolor, mathrsfs, enumitem, setspace, yfonts; the wrapper is the lane's pinned \texttt{amsart} editorial wrapper at 10pt on A4, which restates the theorem-like environments the retained text needs and adds the article's own macro definitions that the retained text needs (\texttt{\textbackslash norm}), with extra wrapper packages (bm, bbm, dsfont, xspace, cancel, mathtools). The delivered numbering is the unit's own. Assets: no retained span contains a figure environment, a graphics-inclusion command, a PDF-page inclusion, a table or a TikZ picture; no image file is copied into this package, the package declares no asset dependency and no image file is redistributed with it. Licence: version arXiv:2508.12654v1 of this article is distributed under the Creative Commons Attribution 4.0 International licence, as recorded in the arXiv licence metadata for this exact version and shown on the abstract page https://arxiv.org/abs/2508.12654v1. A full rights notice is delivered at licenses/arxiv-2508.12654-CC-BY-4.0.txt. Characters: every retained excerpt is pure ASCII, so the delivered engine, XeLaTeX with its default Latin Modern text font, reports no missing character.
\begin{lemma}\label{local structure}
    Let $ p\,\in\, S^{m}(Y)$ be given by $ p\,=\, \sum_{i=1}^{k}a_{i}y_{i},$ where $y_{i} $' s are distinct, and $ a_{1}\,\geq\, a_{2}\,\geq\, \cdots\, \geq\, a_{k}.$  Let $ \pi\,=\,(a_{1},a_{2},\cdots,a_{k})$. Then there is an isomorphism of germs 
    $\prod_{i=1}^{k} X_{a_{i}}^g\,\xrightarrow{\,\,\,\phi\,\,\,}\, \langle S^{m}Y, p \rangle$ such that the following holds:
    
    For any $\pi^{\prime}\,=\, (b_{1}, b_{2},\cdots,b_{\ell})$ such that $ \pi\, \geq\, \pi^{\prime}.$ Then we have $$ \phi^{-1}\langle W_{\pi^{\prime},m}(Y), p\rangle\, =\, \bigcup_{\pi^{\prime}\xrightarrow{r} \pi} \prod_{i=1}^{k} W_{r_{(i)},a_{i}}^{g}.$$ Here, $r$ runs over all refinements of $ \pi^{\prime}$ by $\pi,$ up to equivalence, and $r_{(i)}$ denotes the partition $ (b_{j})_{j\in r^{-1}(i)}$ of $a_i.$


\end{lemma}

\begin{proof}
First, we define the isomorphism $\phi$.  Since $Y$ is smooth, for each $y\in Y$ we have an isomorphism of germs: 
$$ \langle T_{y}Y,0 \rangle\, \xlongrightarrow{\,\cong\,}\,\langle Y, y \rangle .$$ For each $ y\,\in\, Y$, denote this isomorphism by $$ v\, \longmapsto\, y\,\oplus\, v.$$

Let $ \underline{q}\,=\, (\underbrace{y_{1},\cdots,y_{1}}_{a_{1}\text{-times}}, \underbrace{y_{2},\cdots,y_{2}}_{a_{2}\text{-times}},\cdots, \underbrace{y_{k},\cdots,y_{k})}_{a_{k}\text{-times}}\,\in\, Y^{m}.$ Under the action of $S_{m}$ on $Y^{m}$, $ H\,:=\,\textnormal{stab}(\underline{q})\,=\,S_{a_{1}} \,\times\, S_{a_{2}}\,\times\, \cdots\, \times\, S_{a_{k}},$ as $y_{i} \hspace{1pt}$'s are distinct. 
We have $ \langle Y^{m}, \underline{q} \rangle \,\cong\,\underset{i} {\prod} \langle Y^{a_{i}}, (y_{i},y_{i}, \cdots, y_{i})\rangle \,\cong\, \underset{i} {\prod}\langle(T_{y_{i}}Y)^{a_{i}}, \underline{0}\rangle $, and the action of $ H\,=\, \underset{i} {\prod}S_{a_{i}}$ on $ \underset{i} {\prod} \langle (T_{y_{i}}Y)^{a_{i}}, \underline{0}\rangle$ is just the product action. So, we have isomorphisms

$$\hspace{-1.3cm} \langle S^{m}Y,p \rangle\, \cong\, \langle Y^{m}, \underline{q} \rangle \big{/}H$$

    \hspace{185pt}$ \cong\, \underset{i} {\prod} \langle (T_{y_{i}}Y)^{a_{i}}, \underline{0}\rangle\big{/}S_{a_{i}}$

    $ \hspace{185pt}\cong\, \underset{i} {\prod}\langle S^{a_{i}}\mathbb{A}^n,0 \rangle$

    $ \hspace{185pt}\cong\, \underset{i} {\prod}X_{a_{i}}^{g}$.
    
Call this composition of isomorphisms $\phi\,:\,\prod_{i=1}^{k} X_{a_{i}}^g\longrightarrow \langle S^{m}Y, p \rangle$.
So, $ \phi$ is given by 
\begin{equation}\label{phi}
    \phi ((v_{i1}+ \cdots + v_{ia_{i}})_{i})\,=\, \sum_{i}(y_{i}\,\oplus\, v_{i1}+ \cdots + y_{i}\,\oplus\, v_{ia_{i}}).
\end{equation} 
Here, $v_{ij}\in T_{y_{i}}Y,$ and $ v_{i1}+\cdots +v_{ia_{i}}$ is considered as a zero cycle in $ T_{y_{i}}Y$ (it is not the vector addition).





Now we show that $\phi$ has the desired property.
Choose metrics on each $T_{y_{i}}Y$. Let $\epsilon \,>\,0$ be such that $ y_{i}\,\oplus\, v$ is defined for all $i$ and $\norm{v}\, <\, \epsilon, $ and $ y_{i}\,\oplus\, B_{\epsilon}(0)$' s are disjoint. Let $$\widehat{X_{a_{i}}^{g}}\,=\, \Big \{ \sum_{j=1}^{a_{i}}v_{ij} \in S^{a_{i}}(T_{y_{i}}Y) \hspace{1pt}\,\big{|}\, \hspace{1pt}  \norm{v_{ij}}\,<\, \epsilon  \hspace{0.3cm} \textnormal{for all}\,\, j \Big \}.$$ 
So, $\phi$ is represented by a holomorphic open map 
    $\underset{i} {\prod}\widehat{X_{a_{i}}^{g}}\, \xlongrightarrow{\,\,\,\widehat{\phi}\,\,\,}\, S^{m}Y$, which is given by the same formula as in  \eqref{phi}, and is an isomorphism onto its image, which we can call $ \widehat{S^{m}Y}.$ Let $ \widehat{W}_{r_{(i)},a_{i}}^{g}\,=\, \widehat{X_{a_{i}}^{g}} \,\cap\, W_{r_{(i)},a_{i}}^{g} $. It suffices to show $$ {\widehat{\phi}}^{-1}(W_{\pi^{\prime},m}(Y))\,=\, \bigcup_{\pi^{\prime}\xrightarrow{r} \pi} \prod_{i=1}^{k} \widehat{W}_{r_{(i)},a_{i}}^{g} .$$
     \underline{$ {\widehat{\phi}}^{-1}( W_{\pi^{\prime},m}(Y)) \supseteq\underset{\pi^{\prime}\xrightarrow{r} \pi}{\bigcup} \prod_{i=1}^{k} \widehat{W}_{r_{(i)},a_{i}}^{g}  $:} 
    
    Let  $ \pi^{\prime}\,=\, (b_{s})_{s},$ and $ \pi^{\prime}\xrightarrow{r} \pi$ be a refinement. Let $ \underline{v}\,=\,(\underset{s\in r^{-1}(i)}{\sum}b_{s}v_{is})_{i}$ be an element of $ \prod _{i=1}^k \widehat{W}_{r_{(i)},a_{i}}^{g}$, where $ v_{is}\in T_{y_{i}}Y.$ We have 
    $$ \widehat{\phi}(\underline{v})\,=\, \sum_{i} \sum_{s\in r^{-1}(i)}b_{s}y_{i}\, \oplus\, v_{is}\,=\, \sum_{s}b_{s}y_{r(s)}\,\oplus\, v_{r(s),s} \in W_{\pi^{\prime},m}(Y).$$
\underline{$ \widehat{\phi}^{-1}( W_{\pi^{\prime},m}(Y)) \,\subseteq\, \underset{\pi^{\prime}\xrightarrow{r} \pi}{\bigcup} \prod_{i=1}^{k} \widehat{W}_{r_{(i)},a_{i}}^{g}  $:}

    Let $ v_{ij} \in T_{y_{i}}Y$, for $1\,\leq\, i\,\leq\, k,\quad 1\,\leq\, j\, \leq\, a_i,$ be such that $ \underline{v}\,=\, (\sum_{j}v_{ij})_{i} \in \widehat{\phi}^{-1}(W_{\pi^{\prime},m}(Y)),$ i.e., $$ \sum_{i}\sum_{j}y_{i}\,\oplus\, v_{ij} \in W_{\pi^{\prime},m}(Y).$$ So, there are $z_{s}$'s in $Y$ such that 
    $$ \sum_{i=1}^{k} \sum_{j=1}^{a_{i}} y_{i}\,\oplus\, v_{ij}\,=\,\sum_{s=1}^l b_{s}z_{s}.$$
    Since  $ y_{i}\,\oplus\, B_{\epsilon}(0)$'s are disjoint, for each $s$ we have unique $ r(s)\in [k]$ such that $ z_{s} \in y_{r({s})}\, \oplus\, B_{\epsilon}(0).$ This gives a map $ r : [\ell]\longrightarrow [k],$ we have 
    
    \begin{equation}\label{sum}
        \sum_{i=1}^{k}\sum_{j=1}^{a_{i}} y_{i}\,\oplus\, v_{ij}\,=\, \sum_{i=1}^{k}\sum_{s\in r^{-1}(i)}b_{s}z_{s}.
    \end{equation}
    So, $$ \sum_{i=1}^k(\sum_{j=1}^{a_{i}} y_{i}\,\oplus\, v_{ij}\,-\, \sum_{s\in r^{-1}(i)}b_{s}z_{s})\,=\,0.$$
    The $0$-cycle $\sum_{j=1}^{a_{i}} y_{i}\,\oplus\, v_{ij}\,-\, \sum_{s\in r^{-1}(i)}b_{s}z_{s}$ is supported in $y_{i}\,\oplus\, B_{\epsilon}(0).$
    Since $ y_{i}\,\oplus\, B_{\epsilon}(0)$'s are disjoint, we must have 
    \begin{equation}\label{exponential}
\sum_{j=1}^{a_{i}}y_{i}\,\oplus\, v_{ij}\,=\, \sum_{s\in r^{-1}(i)}b_{s}z_{s}\quad \textnormal{for all} \hspace{3pt}i. 
\end{equation}

Taking degrees, we get $ a_{i}\,=\,\underset{s\in r^{-1}(i)}{\sum}b_{s}.$ Thus, $ r$ is a refinement. Since $$ B_{\epsilon}(0)\, \longrightarrow\, y_{i}\,\oplus\, B_{\epsilon}(0)$$
$$ v \longmapsto y_{i}\,\oplus\, v$$
is an isomorphism,  there are unique $ w_{s}\,\in\, B_{\epsilon}(0)$ such that $ z_{s}\,=\, y_{r(s)} \,\oplus\, w_{s}.$ Now, \eqref{exponential} gives $ \sum_{j=1}^{a_{i}} v_{ij}\,=\, \sum _{s\in r^{-1}(i)} b_{s}w_{s}.$ So, $ \sum_{j}v_{ij}\, \in\, \widehat{W}_{r_{(i)},a_{i}}^{g}$. Hence $ \underline{v}\, \in\, \prod_{i=1}^{k}\widehat{W}_{r_{(i)},a_{i}}^{g}$.
\end{proof}

\end{document}
